\section{Data}
\label{sec:data}

The empirical analysis uses a community-district-by-year panel of housing production in New York City.
The panel covers annual housing production from 1970 through 2025, so the same neighborhood-scale units can be followed before and after the court-induced shift from borough-scale to district-scale land-use review.

\subsection{Community-District Panel}

The sample contains New York City's 59 standard community districts.
Community districts are not the political units through which member deference operates, but they are stable neighborhood geographies that can be observed consistently across the sample period.
The design asks whether housing production shifted away from homeowner-heavy neighborhoods as land-use authority became more localized.

\subsection{Homeownership Exposure}

The treatment variable measures how homeowner-heavy a community district was at baseline relative to other districts in the same borough.
I use owner-occupied and occupied housing-unit counts from the Department of City Planning's 1990 community-district profiles.
The main treatment is this rate standardized within borough.
A value of zero means the district had average homeownership for its borough; positive values indicate a more homeowner-heavy district; negative values indicate a more renter-heavy district.

The within-borough scale is important for interpretation.
Homeownership, density, housing stock, and development pressure differ sharply across boroughs.
The comparison is therefore not Staten Island against Manhattan, but homeowner-heavy districts against lower-homeownership districts within the same borough.

\subsection{Housing Production}

Housing production is measured from 25v4 MapPLUTO.
I assign lots to community districts and use the \texttt{yearbuilt} field and residential-unit counts to construct district-year measures of new housing.
The main outcomes are total residential units and units on building-size margins where discretionary review should matter differently.
The descriptive figures show total units, small-building units, and very large buildings.
The main figures use raw unit counts throughout.
The event-study results focus on raw units in 5+ unit buildings and plot the 1--4 unit margin as a placebo.

MapPLUTO has one important limitation.
It is a surviving-stock year-built proxy, not a clean annual completion-flow series.
That limitation becomes more important further back in time, because older buildings are only observed if they survive into the current MapPLUTO vintage.
The current results are therefore descriptive and best interpreted as evidence on whether post-reform production shifted across neighborhoods, rather than as a complete measure of every unit completed in each year.
